\chapter{System przepisywania słów}

W pracy \textit{Active Automata Learning with Advice}~\cite{srs} przedstawiono podejście, w którym algorytm aktywnego uczenia się 
może otrzymać dodatkową podpowiedź, dzięki której możliwe jest zredukowanie liczby zapytań kierowanych do nauczyciela.
Autorzy sugerują także w jaki sposób wykorzystać to podejście w innych algorytmach aktywnego uczenia się automatów,
z wykorzystaniem odpowiednika Twierdzenia Myhilla-Nerode'a.

Podejście to nie wymaga zmian w algorytmie ucznia. Zapytania, które zadaje uczeń nie są wprost przekazywane do nauczyciela,
tylko kierowane do systemu, który wnioskuje o nich na podstawie zadanego systemu SRS.
Jeśli system nie jest w stanie odpowiedzieć na zapytanie, wtedy przekierowuje je do nauczyciela i zwraca jego odpowiedź.

Z punktu widzenia implementacji przedstawionej w tej pracy szczególnie istotne jest wnioskowanie
podczas zapytań o równoważność, które stanowią znaczący narzut czasowy.
Zapytania o przynależność nie są kosztowne w przypadku rozważanego algorytmu,
dlatego ta część nie została wykorzystana.

Jan Otop i Michał Fica wprowadzają następujące definicje: system SRS $R$ jest spójny z językiem, wtedy i tylko wtedy,
gdy dla wszystkich słów $s, t$, jeśli $s \to^*_R t$, to $s \in L \Leftrightarrow t \in L$.
Podpowiedzią dla algorytmu jest system SRS $R$ spójny z językiem docelowym.

Następnie dowodzą, że jeśli $R$ jest SRS nad $\Sigma$ i $A$ jest minimalnym DFA również nad $\Sigma$,
to SRS $R$ jest spójny z $L(A)$ wtedy i tylko wtedy, gdy dla każdego stanu $q \in Q_A$ i każdej reguły
$l \to r \in R$ mamy $\delta_A(q, l) = \delta_A(q, r)$.

Ta obserwacja umożliwia przeiterowanie się po konfiguracjach automatu DFA i regułach systemu SRS $R$ w celu wykrycia niespójności hipotezy z systemem.
Następnie pokazują jak znając stan $q \in Q_A$ i $l \to r \in R$, takie że tylko jeden ze stanów $H(q, l)$ i $H(q, r)$ należy do $F_H$
skonstruować kontrprzykład przy użyciu dokładnie jednego zapytania $MQ$.

\textbf{Problemy z adaptacją tego podejścia.}

\begin{enumerate}
    \item Zastosowany odpowiednik Twierdzenia Myhilla-Nerode'a dla automatów visibly pushdown nie jest prawo-kongruencją. 
    Oznacza to, że dwie różne konfiguracje hipotezy $H$ mogą generować dokładnie te same prawostronne języki.
    Dlatego znalezienie reguły $r \to l \in R$ i konfiguracji $(q, s)$, takich że konfiguracje $H(q, l, s)$ i $H(q, r, s)$
    się różnią nie gwarantuje nam, że automat nie jest zgodny z systemem SRS $R$.
    
    \item Nie jest możliwe sprawdzenie wszystkich osiągalnych konfiguracji dla hipotezy, bo w ogólności jest ich nieskończenie wiele.
    Musimy ograniczyć się tylko do części osiągalnych konfiguracji.
\end{enumerate}

Z powodu trudności w zastosowaniu tego podejścia, zostały nałożone dodatkowe obostrzenia na system SRS, który może być podpowiedzią dla algorytmu.
Każda reguła postaci $l \rightarrow r$ powinna składać się ze słów $l$ oraz $r$, które są dobrze zbalansowane. Ponieważ nie zależą one od stosu,
ani nie zmieniają jego zawartości, odpowiadają one pewnemu przekształceniu stanów analogicznie jak dla automatów DFA.

Umożliwia to sprawdzenie czy automat jest zgodny z daną regułą $l \rightarrow r$, poprzez sprawdzenie równoważności dwóch automatów $H_l$ oraz $H_r$, utworzonych
na podstawie hipotezy $H$, z alfabetem rozszerzonym o dodatkową literę $l'$ typu local.
Automat $H_l$ po przeczytaniu litery $l'$ trafia do stanu, w którym znalazłby się automat $H$ po przeczytaniu słowa $l$, a $H_r$ po przeczytaniu $l'$
trafia do stanu w którym znalazłaby się hipoteza po przeczytaniu słowa $r$. Jest to możliwe, gdyż słowa dobrze
zbalansowane nie zależą od stosu ani go nie modyfikują.

Problemem tego podejścia wydawało się sprawdzanie równoważności, gdyż w~aktualnej implementacji jest wykorzystywany ten sam algorytm, który sprawdza
równoważność hipotezy z automatem docelowym. Jednak hipoteza przez znaczący czas trwania algorytmu jest znacznie mniejsza niż automat docelowy, przez co produkt
$H_l \times H_r$ jest znacząco mniejszy niż $H \times T$, gdzie $T$ jest automatem docelowym. W praktyce sprawia to,
że sprawdzanie zgodności z systemem SRS może być znacznie tańsze niż sprawdzanie równoważności hipotezy z automatem docelowym.

\textbf{Finalne zastosowanie} rozszerza ideę systemu o parametryzowane reguły postaci $l_1Xl_2 \rightarrow r_1Xr_2$, reguła ta reprezentuje tożsamość
$wl_1Xl_2v \in L \Leftrightarrow wr_1Xr_2v \in L$, gdzie $w,v \in \Sigma^*$, $X \in \Sigma^*$ jest słowem dobrze zbalansowanym oraz słowa $l_1l_2$ i $r_1r_2$
są dobrze zbalansowane.

Algorytm w celu sprawdzenia zgodności hipotezy $H$ z systemem SRS $R$ rozpoczyna od wyznaczenia zbioru 
słów dobrze zbalansowanych, które mogą zostać podstawione za parametr $X$.
Ponieważ liczba tych słów powinna być ograniczona, obliczana jest wartość każdego słowa
na podstawie której podejmowana jest decyzja, czy słowo to powinno zostać wykorzystane.
Dystansem $dist(w_1, w_2)$ pomiędzy słowami $w_1$ i $w_2$ określamy liczbę takich stanów
$q \in Q$, że $H(q, \bot, w_1) \neq H(q, \bot, w_2)$.
Następnie w momencie dodawania słowa do zbioru słów dobrze zbalansowanych określana jest jego wartość,
jako minimalny dystans od już zawartych w tym zbiorze słów. Jeśli zbiór ten zawiera zbyt wiele elementów
usuwane są słowa o najmniejszej wartości. Podejście to ma na celu zmaksymalizować szanse algorytmu na
wykrycie niespójności.

Następnie dla każdej reguły oraz słowa dodaje nową literę lokalną do automatów $H_l$ oraz $H_r$,
po czym sprawdzana jest ich równoważność. Jeżeli istnieje słowo $w$, które należy do dokładnie jednego
z języków generowanych przez automaty $H_l$ oraz $H_r$ to oznacza, że została wykryta niespójność
z systemem SRS. Należy wtedy skonstruować słowa $w_l$ oraz $w_r$ poprzez podstawienie za wszystkie dodatkowe
litery lokalne odpowiednio lewej oraz prawej strony reguły SRS odpowiadającej danej literze.
Ostatecznie jedno z tych słów jest poprawnym kontrprzykładem, ponieważ system SRS jest zgodny z automatem docelowym.
Umożliwia to podjęcie decyzji, które słowo należy zwrócić, przy pomocy jednego zapytania \textbf{MQ}.

Podejście to gwarantuje, że w przypadku wykrycia niespójności hipoteza $H$ nie jest zgodna z systemem SRS.
Nie gwarantuje jednak wykrycia każdej możliwej niezgodności hipotezy z tym systemem.